\section{Conclusion}

We presented \textsc{ReasoningGuard}, a dual-layer defense framework for MCP-based LLM agents that combines Protocol-Attested Tool Gateway (PTG) with Reasoning Trace Verifier (RTV). Grounded in the observability-based defense limitation property, our approach addresses both protocol-level (L4) and cognitive-level (L2) attacks that no single-layer defense can comprehensively cover. PTG extends AttestMCP with semantic intent attestation, providing an intent-aware capability boundary that ablation identifies as the most critical single defense component. RTV provides evidence consistency checking with three anomaly classes (OAV, IAD, CAI), achieving 100\% TPR and 0\% FPR on direct evaluation and providing clear incremental value beyond PTG on L2 attacks. We position RTV as an evidence consistency checker rather than a reasoning correctness verifier, acknowledging that LLM-generated traces are post-hoc justifications that may not faithfully reflect internal computation. Across eight SOTA models and 9{,}192 evaluation episodes, \textsc{ReasoningGuard} consistently achieves the lowest ASR, reducing it by 48.5\%--92.6\% relative to No Defense with 0.1\,ms defense overhead and a measurable TCR cost (16.7\,pp on GPT-4o). We identify three adaptive adversary evasion strategies and discuss their feasibility. Our MCPTox+ benchmark of 342 scenarios across six attack categories provides the first comprehensive evaluation of cross-session (T3) memory poisoning in MCP security. Future work should explore stronger detection of social engineering injections, neural judge models for adaptive adversary resistance, formal authorization integration, and extension to multi-agent architectures.
